The third stage of the project was the frontend.
This was a vital stage is actually producing a functioning visualisation
that would later be backed by the relational database content.
Four sections were outlined for this stage: JavaScript vs TypeScript,
ESLint \&Prettier, developer relationship algorithm, \&applying D3.js.

The first choice that had to be made was what programming language should be used for the project.
The two candidates for this were: JavaScript and TypeScript.
Firstly, JavaScript was and will always be incredibly well-known in the frontend development community,
making it the common go-to option.
However, the lack of object-oriented programming concepts and tooling support makes it less enticing.
For instance, being able to detect errors early on in a large project will be beneficial in the long-term,
which is one of the many features that TypeScript has to offer over JavaScript.
In contrast, TypeScript is purely statically typed which enhances error detection and, thus, ultimately, code quality.
Additionally, TypeScript supports the latest ECMAScript features whilst also be a direct subsidiary of JavaScript,
meaning that it supports all of its features by default.
The primary downside to choosing TypeScript is the learning curve for amateur developers
who are not used to typing statically.
There are also extra compilation steps that are required to fully implement this static typing feature,
therefore creating extra overhead as well.
Regardless of this,
TypeScript was ultimately chosen as the frontend programming language for this project due to its cutting-edge nature and abundance of features superior to JavaScript.

ESLint and Prettier were invaluable tools for the frontend development of this project.
ESLint is a styling tool that, through static typing analysis,
reports any compilation errors before the compilation process actually begins.
For this project,
the airbnb and airbnb-typescript ESLint configurations were used
to ensure all the project code matched industry-level standards.
All the necessary rules are stored inside a .eslintrc.cjs file.

The complete ESLint configuration file is shown below:

\begin{verbatim}
    module.exports = {
      root: true,
      env: {browser: true, es2020: true},
      extends: [
        'eslint:recommended',
        'plugin:@typescript-eslint/recommended',
        'plugin:react-hooks/recommended',
        'airbnb',
        'airbnb-typescript',
        'plugin:react/jsx-runtime',
      ],
      ignorePatterns: ['dist', '.eslintrc.cjs', 'vite.config.ts'],
      parser: '@typescript-eslint/parser',
      parserOptions: {
        project: './tsconfig.json',
      },
      plugins: ['react-refresh'],
      rules: {
        'react-refresh/only-export-components': [
          'warn',
          {allowConstantExport: true},
        ],
        'linebreak-style': [
          'error',
          (process.platform === 'win32' ? 'windows' : 'unix')
        ]
      },
    }
\end{verbatim}

Additional attributes such as extends, parserOptions,
and linebreak-style are used
to customise the configuration to work correctly with all the project's installed dependencies.

On top of that, Prettier is also used.
This tool focuses more on the formatting of the project code rather than the styling of it.
Thus, the readability of the code is more accessible and the project is easier to contribute to in industry.
A bonus of using Prettier is that, unlike ESLint,
it does not require a separate configuration file to function correctly.
In fact, the project just used to default Prettier settings without a separate configuration being necessary.

A crucial part of the frontend stage was establishing
how the visualisation would associate development with their contributions to a project's pull requests.
The visualisation worked by determining the size of a node
(i.e.\ a developer) by how many comments they made on the individual project.
The links between nodes (i.e.\ other developers)
were created when one developer responded to another developer's comment on the SAME pull request.
The thickness of these links between nodes (i.e.\ other developers)
was determined by how many comments two developments made between each other on the SAME pull request.

With that context in mind, the relationship between developers was determined in the following way:

\begin{enumerate}
    \item The system begins by initialising a list of pull requests, nodes, and links (to structure the graph).
    \item
    All the pull requests (that are stored within the relational database) are stored in the pull request list.
    \item For each pull request, if the author of the pull request is not already a node that exists in the graph (i.e.\ not a member of the nodes list), then add it with a size of zero.
    \item For each comment in each pull request, if the author of the comment is not already a node that exists in the graph (i.e.\ not a member of the nodes list), then add it with a size of zero.
    \item If they do exist, then increment the size of that node by one.
    \item Store the author of the pull request and the author of the comment in a list, and sort this list.
    \item If this pair of users (in the list) is not already in the links list, then add them with a thickness of zero.
    \item Otherwise, increment the thickness between the pair of users by one.
\end{enumerate}

The final step of the frontend milestone was to apply D3.js to have a visualisation that users can actually view.
For the use case of the project, a force-directed graph was chosen as the appropriate D3 component.
This was due to the easy ability to add nodes and links in a JSON format that could be easily retrieved from the backend API endpoint.
An important point to mention is that the backend API endpoint was only callable thanks to the axios npm library.
Axios allowed the API endpoint to be called and receive and store the response in a single variable, with a ``data'' attribute being used to retrieve it.
The previously mentioned developer relationship algorithm was what provided the data compatible with the force-directed graph.
In particular, this required a JSON response with ``nodes'' and ``links'' objects, respectively.
These objects also required their own specific attributes, such as size for ``nodes'' and thickness for ``links''.
Several D3 forces and physics-related components were applied to make the visualisation successful due to the structure of the graph itself.
With some sizes and thicknesses being inconsistent with each other (depending on the time period scale that was chosen), making sure the physics of the graph was smooth was important.
It took some testing and trial-and-error to ensure that the physics were professional-looking but also practical, but eventually everything came into place.
One major benefit to using D3 was how it seamlessly applied normalisation and scaling to quantitative scales.
For instance, a separate function was required to normalise and scale the thickness of each link between nodes since at first they were too varying in size.
This was needed in other numerical aspects of the graph as well.
Another surprise when finally bringing the frontend together was how efficient the performance was when all the nodes were added to the graph.
Despite load times reaching upwards of one minute, the graph component was stable enough to be usable despite thousands of nodes being present.